% GENERATED FILE. Edit canonical lessons, then run npm run generate:books.
% canonical-source-hash: fnv1a64:50a2b45d2d9f2ba9
% canonical-lessons: JA-C336-kankyou, JA-C336-shizen, JA-C336-chikyuu, JA-C336-kuuki, JA-C336-osen

\chapter{The Environment}
\label{ch:ja-the-environment}
\hlchaptermodality{336}

\begin{chapteropening}
\textbf{By the end of this chapter:} \emph{I can give my opinion about the environment and say why.}

Complete the last lesson of chapter 336: I can give my opinion about the environment and say why.
\end{chapteropening}

\section[\texorpdfstring{\ja{かんきょう} (kankyō)}{かんきょう (kankyō)}]{\ja{かんきょう} (kankyō) — the environment; surroundings}

\label{lesson:JA-C336-kankyou}



\begin{quote}
Before the new one: say the Japanese for oneself, then the Japanese for other people.
\end{quote}

\subsection*{You'll want to know: \ja{かんきょう}}
\textbf{\ja{かんきょう}} — \emph{kankyō} — \textquotedblleft{}the environment; surroundings\textquotedblright{}.

\begin{quote}
\textbf{\ja{かんきょうの} \ja{ために、じてんしゃを} \ja{つかいます。}} — \emph{kankyō no tame ni, jitensha o tsukaimasu} — \textquotedblleft{}I ride a bicycle for the sake of the environment.\textquotedblright{}
\end{quote}

\begin{cousinweb}
\textbf{\ja{かんきょう}} (\emph{kankyō}) is Sino-Japanese (kango): two Chinese characters read with their on'yomi, readings borrowed from Chinese. They are \emph{kan}, \textquotedblleft{}ring\textquotedblright{} (Middle Chinese \emph{hwaen}) and \emph{kyō}, \textquotedblleft{}boundary\textquotedblright{} (KANJIDIC2; Middle Chinese in Baxter's transcription).
\end{cousinweb}

\begin{grammarlens}[title={What you've built}]
One more word for giving reasons and opinions.
\end{grammarlens}

\subsection*{Your turn}
\begin{itemize}
\raggedright
  \item \emph{Say it:} \emph{kankyō}
  \item \emph{Say it:} \emph{kankyō}, once more
  \item \emph{Say it:} \emph{kankyō no tame ni, jitensha o tsukaimasu} — I ride a bicycle for the sake of the environment.
  \item \emph{Recall:} say the Japanese for oneself, then the Japanese for other people, then say \emph{kankyō} again
\end{itemize}

\begin{tcolorbox}[breakable,colback=teal!4,colframe=teal!35!black,title={Before you move on}]
What does \ja{かんきょう} mean? (\textquotedblleft{}The environment; surroundings\textquotedblright{}.) Say it once more.
\end{tcolorbox}
\section[\texorpdfstring{\ja{しぜん} (shizen)}{しぜん (shizen)}]{\ja{しぜん} (shizen) — nature}

\label{lesson:JA-C336-shizen}



\begin{quote}
Before the new one: say the Japanese for other people, then the Japanese for the environment.
\end{quote}

\subsection*{You'll want to know: \ja{しぜん}}
\textbf{\ja{しぜん}} — \emph{shizen} — \textquotedblleft{}nature\textquotedblright{}.

\begin{quote}
\textbf{\ja{いなかは} \ja{しぜんが} \ja{おおいから、すきです。}} — \emph{inaka wa shizen ga ōi kara, suki desu} — \textquotedblleft{}I like the countryside because there is so much nature.\textquotedblright{}
\end{quote}

\begin{cousinweb}
\textbf{\ja{しぜん}} (\emph{shizen}) is Sino-Japanese (kango): two Chinese characters read with their on'yomi, readings borrowed from Chinese. They are \emph{shi}, \textquotedblleft{}oneself\textquotedblright{} (Middle Chinese \emph{dzijH}) and \emph{zen}, \textquotedblleft{}so\textquotedblright{} (\emph{nyen}) (KANJIDIC2; Middle Chinese in Baxter's transcription).
\end{cousinweb}

\begin{grammarlens}[title={What you've built}]
One more word for giving reasons and opinions.
\end{grammarlens}

\subsection*{Your turn}
\begin{itemize}
\raggedright
  \item \emph{Say it:} \emph{shizen}
  \item \emph{Say it:} \emph{shizen}, once more
  \item \emph{Say it:} \emph{inaka wa shizen ga ōi kara, suki desu} — I like the countryside because there is so much nature.
  \item \emph{Recall:} say the Japanese for other people, then the Japanese for the environment, then say \emph{shizen} again
\end{itemize}

\begin{tcolorbox}[breakable,colback=teal!4,colframe=teal!35!black,title={Before you move on}]
What does \ja{しぜん} mean? (\textquotedblleft{}Nature\textquotedblright{}.) Say it once more.
\end{tcolorbox}
\section[\texorpdfstring{\ja{ちきゅう} (chikyū)}{ちきゅう (chikyū)}]{\ja{ちきゅう} (chikyū) — the earth; the globe}

\label{lesson:JA-C336-chikyuu}



\begin{quote}
Before the new one: say the Japanese for the environment, then the Japanese for nature.
\end{quote}

\subsection*{You'll want to know: \ja{ちきゅう}}
\textbf{\ja{ちきゅう}} — \emph{chikyū} — \textquotedblleft{}the earth; the globe\textquotedblright{}.

\begin{quote}
\textbf{\ja{ちきゅうが} \ja{あつく} \ja{なったと} \ja{おもいます。}} — \emph{chikyū ga atsuku natta to omoimasu} — \textquotedblleft{}I think the earth has got hotter.\textquotedblright{}
\end{quote}

\begin{cousinweb}
\textbf{\ja{ちきゅう}} (\emph{chikyū}) is Sino-Japanese (kango): two Chinese characters read with their on'yomi, readings borrowed from Chinese. They are \emph{chi}, \textquotedblleft{}ground\textquotedblright{} (Middle Chinese \emph{dijH}) and \emph{kyū}, \textquotedblleft{}sphere\textquotedblright{} (\emph{gjuw}) (KANJIDIC2; Middle Chinese in Baxter's transcription).
\end{cousinweb}

\begin{grammarlens}[title={What you've built}]
One more word for giving reasons and opinions.
\end{grammarlens}

\subsection*{Your turn}
\begin{itemize}
\raggedright
  \item \emph{Say it:} \emph{chikyū}
  \item \emph{Say it:} \emph{chikyū}, once more
  \item \emph{Say it:} \emph{chikyū ga atsuku natta to omoimasu} — I think the earth has got hotter.
  \item \emph{Recall:} say the Japanese for the environment, then the Japanese for nature, then say \emph{chikyū} again
\end{itemize}

\begin{tcolorbox}[breakable,colback=teal!4,colframe=teal!35!black,title={Before you move on}]
What does \ja{ちきゅう} mean? (\textquotedblleft{}The earth; the globe\textquotedblright{}.) Say it once more.
\end{tcolorbox}
\section[\texorpdfstring{\ja{くうき} (kūki)}{くうき (kūki)}]{\ja{くうき} (kūki) — air; the atmosphere}

\label{lesson:JA-C336-kuuki}



\begin{quote}
Before the new one: say the Japanese for nature, then the Japanese for the earth.
\end{quote}

\subsection*{You'll want to know: \ja{くうき}}
\textbf{\ja{くうき}} — \emph{kūki} — \textquotedblleft{}air; the atmosphere\textquotedblright{}.

\begin{quote}
\textbf{\ja{まちは} \ja{くうきが} \ja{きたないので、いなかに} \ja{すみたいです。}} — \emph{machi wa kūki ga kitanai node, inaka ni sumitai desu} — \textquotedblleft{}The air in town is dirty, so I want to live in the country.\textquotedblright{}
\end{quote}

\begin{cousinweb}
\textbf{\ja{くうき}} (\emph{kūki}) is Sino-Japanese (kango): two Chinese characters read with their on'yomi, readings borrowed from Chinese. They are \emph{kū}, \textquotedblleft{}empty\textquotedblright{} (Middle Chinese \emph{khuwng}) and \emph{ki}, \textquotedblleft{}air\textquotedblright{} (\emph{khj+jH}) (KANJIDIC2; Middle Chinese in Baxter's transcription).
\end{cousinweb}

\begin{grammarlens}[title={What you've built}]
One more word for giving reasons and opinions.
\end{grammarlens}

\subsection*{Your turn}
\begin{itemize}
\raggedright
  \item \emph{Say it:} \emph{kūki}
  \item \emph{Say it:} \emph{kūki}, once more
  \item \emph{Say it:} \emph{machi wa kūki ga kitanai node, inaka ni sumitai desu} — The air in town is dirty, so I want to live in the country.
  \item \emph{Recall:} say the Japanese for nature, then the Japanese for the earth, then say \emph{kūki} again
\end{itemize}

\begin{tcolorbox}[breakable,colback=teal!4,colframe=teal!35!black,title={Before you move on}]
What does \ja{くうき} mean? (\textquotedblleft{}Air; the atmosphere\textquotedblright{}.) Say it once more.
\end{tcolorbox}
\section[\texorpdfstring{\ja{おせん} (osen)}{おせん (osen)}]{\ja{おせん} (osen) — pollution}

\label{lesson:JA-C336-osen}



\begin{quote}
Before the new one: say the Japanese for the earth, then the Japanese for air.
\end{quote}

\subsection*{You'll want to know: \ja{おせん}}
\textbf{\ja{おせん}} — \emph{osen} — \textquotedblleft{}pollution\textquotedblright{}.

\begin{quote}
\textbf{\ja{かわの} \ja{おせんが} \ja{ひどいから、さかなが} \ja{いません。}} — \emph{kawa no osen ga hidoi kara, sakana ga imasen} — \textquotedblleft{}The river is badly polluted, so there are no fish.\textquotedblright{}
\end{quote}

\begin{cousinweb}
\textbf{\ja{おせん}} (\emph{osen}) is Sino-Japanese (kango): two Chinese characters read with their on'yomi, readings borrowed from Chinese. They are \emph{o}, \textquotedblleft{}dirty\textquotedblright{} (Middle Chinese \emph{'u}) and \emph{sen}, \textquotedblleft{}dye\textquotedblright{} (\emph{nyemX}) (KANJIDIC2; Middle Chinese in Baxter's transcription).
\end{cousinweb}

\begin{grammarlens}[title={What you've built}]
One more word for giving reasons and opinions.
\end{grammarlens}

\subsection*{Your turn}
\begin{itemize}
\raggedright
  \item \emph{Say it:} \emph{osen}
  \item \emph{Say it:} \emph{osen}, once more
  \item \emph{Say it:} \emph{kawa no osen ga hidoi kara, sakana ga imasen} — The river is badly polluted, so there are no fish.
  \item \emph{Recall:} say the Japanese for the earth, then the Japanese for air, then say \emph{osen} again
\end{itemize}

\begin{tcolorbox}[breakable,colback=teal!4,colframe=teal!35!black,title={Before you move on}]
What does \ja{おせん} mean? (\textquotedblleft{}Pollution\textquotedblright{}.) Say it once more.
\end{tcolorbox}
